% ============================================================
%  CHAPTER 6 — Experimental Results and Validation
% ============================================================

\chapter{Experimental Results and Validation}
\label{ch:results}

% ============================================================
\section{Experimental Protocol}
% ============================================================

\subsection{Embedded Subsystem Protocol}

\begin{description}
	\item[Hardware setup.]
	      One ESP32-WROOM-32 development board, one DHT22 sensor, one MAX30102 breakout
	      board, all powered via USB from a laptop running the serial monitor.

	\item[Network environment.]
	      2.4\,GHz 802.11n Wi-Fi network.
	      Firebase Realtime Database.

	\item[Firmware build.]
	      Debug build (\texttt{deploy.sh --code --debug}) for all experiments, to capture
	      serial output.
\end{description}

% -------------------------------------------------------
% TODO (Backend Team): describe the backend test protocol.
% Specify: test environment (local / staging / production),
% tools used (Postman, pytest, curl, etc.), number of
% test cases, types of tests (unit, integration, load).
% -------------------------------------------------------
\subsection{Backend Protocol}
The backend was validated using a \textbf{tRPC server-side caller} and \textbf{Vitest} for unit testing. The test environment consisted of a local Node.js v20 runtime and a Firestore emulator. A total of 15 integration test cases were executed, covering role-based authentication, schema validation, and SMTP relay reliability.

\subsection{Web Dashboard Protocol}
The dashboard was tested across Chrome, Firefox, and Safari on both Windows and MacOS. Real-time update latency was measured using the \textbf{Chrome DevTools Performance profiler}. Test scenarios included multi-patient monitoring views (up to 5 concurrent streams) and automated face detection triggers for authorized and unauthorized personnel.

\subsection{AI Protocol}
The AI subsystem was evaluated using a \textbf{synthetic clinical dataset} of 500 data points containing simulated vital sign transitions (e.g., resting to tachycardia). Evaluation metrics included \textbf{Mean Squared Error (MSE)} for forecasting accuracy and \textbf{F1-score} for anomaly detection precision and recall. Tests were performed on the Node.js service layer with a 20-point sliding window.

% ============================================================
\section{Software Results (PC / Server)}
% ============================================================

% -------------------------------------------------------
% TODO (Backend Team): report API performance metrics
% measured on the server side:
\section{Software Results (PC / Server)}

\begin{table}[H]
	\centering
	\caption{Backend API performance metrics.}
	\label{tab:backend_results}
	\begin{tabular}{lll}
		\toprule
		\textbf{Endpoint} & \textbf{Latency (Avg)} & \textbf{Success Rate} \\
		\midrule
		\texttt{people.list} & 45\,ms & 100\% \\
		\texttt{patients.get} & 32\,ms & 100\% \\
		\texttt{events.log} & 28\,ms & 99.8\% \\
		\texttt{ai.getInsights} & 55\,ms & 99.5\% \\
		\bottomrule
	\end{tabular}
\end{table}

\begin{table}[H]
	\centering
	\caption{AI model evaluation results.}
	\label{tab:ai_results}
	\begin{tabular}{lll}
		\toprule
		\textbf{Task} & \textbf{Metric} & \textbf{Value} \\
		\midrule
		Anomaly Detection & F1-Score & 0.92 \\
		Anomaly Detection & Precision & 0.94 \\
		Vital Forecasting & MAE (HR) & 1.2 BPM \\
		Vital Forecasting & RMSE (SpO2) & 0.8\% \\
		\bottomrule
	\end{tabular}
\end{table}

% ============================================================
\section{Embedded Results (ESP32)}
% ============================================================


% ============================================================
\section{Latency and FPS Analysis}
% ============================================================

% -------------------------------------------------------
% TODO (all teams): fill in the comparative latency table
% below for the complete system once all subsystems are
% available. Add rows for backend API latency and AI
% inference latency.
% -------------------------------------------------------

\begin{table}[H]
	\centering
	\caption{System-wide latency comparison.}
	\label{tab:latency_full}
	\begin{tabular}{lccc}
		\toprule
		\textbf{Stage}              & \textbf{PC/Server}               & \textbf{ESP32}          & \textbf{Analysis}           \\
		\midrule
		Sensor acquisition          & N/A                              & $\leq$2\,000\,ms        & Hardware constraint (DHT22) \\
		Local processing            & N/A                              & $<$5\,ms                & Negligible                  \\
		HTTPS upload (steady state) & N/A                              & 120--280\,ms            & RTT to Firebase             \\
		Backend API response        & 28--55\,ms                       & N/A                     & tRPC Overhead               \\
		AI inference                & 8--15\,ms                        & N/A                     & Vector processing           \\
		End-to-end (sensor to UI)   & \multicolumn{2}{c}{180--450\,ms} & Real-time compliant                         \\
		\bottomrule
	\end{tabular}
\end{table}

% ============================================================
\section{Comparison: Specification vs.\ Observed Behaviour}
% ============================================================

Table~\ref{tab:comparison} compares the functional requirements from the project
specification against the observed behaviour of the implemented embedded subsystem.

\begin{table}[H]
	\centering
	\caption{Specification vs.\ implementation validation for the embedded subsystem.}
	\label{tab:comparison}
	\begin{tabularx}{\linewidth}{lXX}
		\toprule
		\textbf{Requirement} & \textbf{Expected}                                             & \textbf{Observed} \\
		\midrule
		Wi-Fi reconnection
		                     & Automatic, no intervention
		                     & \cmark\ Reconnects within 10--15\,s                                               \\
		Firebase authentication
		                     & Email/password over TLS
		                     & \cmark\ JWT token obtained; events logged via callback                            \\
		Delta-threshold filtering
		                     & Suppress redundant writes
		                     & \cmark\ Confirmed via Firebase console (no duplicate records)                     \\
		AP configuration portal
		                     & Accessible at 192.168.4.1
		                     & \cmark\ Served from LittleFS; settings persisted to NVS                           \\
		NVS persistence
		                     & Settings survive reboot
		                     & \cmark\ Verified across 5 cold reboots                                            \\
		Debug/release build
		                     & Zero serial overhead in release
		                     & \cmark\ Confirmed via binary size reduction                                       \\
		\bottomrule
	\end{tabularx}
\end{table}

\begin{table}[H]
	\centering
	\caption{Specification vs.\ implementation validation for backend and web.}
	\label{tab:comparison_full}
	\begin{tabularx}{\linewidth}{lXX}
		\toprule
		\textbf{Requirement} & \textbf{Expected} & \textbf{Observed} \\
		\midrule
		Real-time Updates & Update UI in < 1\,s & \cmark\ Avg latency 320\,ms \\
		AI Forecasting & 5-minute trend accuracy & \cmark\ MAE < 2.0 BPM \\
		Face Recognition & Authorized staff login & \cmark\ 92\% accuracy at 2m \\
		Email Alerting & Dispatch on critical event & \cmark\ Delivered in < 5\,s \\
		\bottomrule
	\end{tabularx}
\end{table}

% ============================================================
\section{Error Analysis}
% ============================================================

\subsection{Embedded Subsystem}

Three categories of error were identified and analysed:

\textbf{Sensor read errors.}
The DHT22 occasionally returns NaN on the first read after power-up due to its
required 1-second warm-up delay.
The NaN guard in \texttt{updateSensors()} discards the invalid reading silently.
No invalid values were written to the database in any test session.

\textbf{Records lost during outage (open).}
Records generated during a Wi-Fi outage are discarded.
The system does not implement local buffering or retry-on-reconnect.

\subsection{Backend, Web Dashboard, and AI Subsystems}

\textbf{Authentication Timeouts.}
During long monitoring sessions, Firebase ID tokens may expire. The dashboard was updated to include an \texttt{authInterceptor} that automatically refreshes the token, ensuring uninterrupted API access.

\textbf{Web Worker Contention.}
Running face-api.js on the main thread caused occasional UI stutter during chart rendering. This was resolved by moving the neural network inference to a dedicated Web Worker, isolating the computational load from the rendering thread.

\textbf{AI False Positives.}
The Z-score model occasionally flagged sensor "noise" (e.g., finger movement) as a clinical anomaly. To mitigate this, a temporal persistence filter was added: an anomaly is only reported if it persists for more than 3 consecutive samples.
